\documentclass[sigconf]{acmart}

\setcopyright{none}
\acmConference[DAI '26]{8th International Conference on Distributed Artificial Intelligence}{November 29--December 2, 2026}{Hong Kong}
\acmBooktitle{8th International Conference on Distributed Artificial Intelligence (DAI '26), November 29--December 2, 2026, Hong Kong}

\title{MetaRoute-Bench: Evaluating Meta-Decision Policies for Agentic Workflows}

% DAI Industry Track is single blind. Replace these placeholders before submission.
\author{Author Names Required}
\affiliation{\institution{Organization Required}\country{Country Required}}
\email{contact@example.com}

\begin{document}

\begin{abstract}
Agentic systems must repeatedly decide whether to answer directly, decompose a task, invoke a tool, execute code, delegate to a specialist, verify an intermediate result, or recover from failure. These meta-decisions affect not only task success but also operating cost and latency, yet they are often embedded inside an orchestration framework and evaluated only through aggregate task accuracy. We present MetaRoute-Bench, an open, inspectable framework for comparing meta-decision policies under a shared execution model. The initial benchmark contains 180 synthetic task profiles spanning data analysis, research, and document processing, eight routing policies, and 30 paired random seeds. Across 43,200 traces, a task-aware compositional policy achieves 79.4\% success compared with 76.7\% for a strong workload-specific static policy, 67.4\% for one-shot task routing, and 52.9\% for direct answering. Relative to the static policy, this is a 2.7 percentage-point improvement (paired 95\% CI: $\pm$2.0 points) at 4.7\% higher mean cost and 6.4\% higher latency. Ablations show the largest losses when route composition is restricted to one operation and when verification is removed. These results are generated by a seeded offline execution model rather than a live deployment; accordingly, our primary contribution is a reproducible evaluation method and an analysis of routing-policy tradeoffs, not evidence of production effectiveness. We release task generation, policies, traces, tests, and analysis artifacts to support live-system validation.
\end{abstract}

\keywords{agentic systems, routing, orchestration, evaluation, tool use, reproducibility}

\maketitle

\section{Introduction}

An agentic application is more than a language model call. It is a distributed decision process involving models, retrieval systems, code executors, specialist agents, verification stages, retry logic, and human escalation. Before any component can help, the system must decide which component to invoke and when. A poor meta-decision can waste time on unnecessary decomposition, call an irrelevant tool, execute unsafe or unhelpful code, or stop before a result has been checked.

Existing research has established the value of interleaving reasoning and action~\cite{yao2023react}, learning API use~\cite{schick2023toolformer,qin2024toolllm}, routing among models~\cite{chen2024frugalgpt,ong2025routellm}, and optimizing function-call plans~\cite{kim2024llmcompiler}. Agent benchmarks reveal persistent failures in long-horizon reasoning and decision making~\cite{liu2024agentbench}, while domain benchmarks increasingly emphasize executable outcomes~\cite{jimenez2024swebench,yao2024taubench}. However, system builders still lack a small, inspectable protocol for isolating the policy that chooses among reasoning modes and measuring its success, cost, and latency tradeoffs.

This paper introduces MetaRoute-Bench, an offline benchmark and trace format for that meta-decision layer. We ask: \emph{under controlled execution assumptions, does task-aware adaptive routing improve operational outcomes over fixed and one-shot policies?} Our contributions are:

\begin{itemize}
  \item a typed framework separating task profiles, routing policies, execution, traces, and evaluation;
  \item a balanced suite of 180 synthetic profiles across three operational workload families;
  \item a paired comparison of eight policies over 30 seeds and 43,200 traces; and
  \item ablations and workload-level analyses identifying when composition, verification, decomposition, and recovery affect outcomes.
\end{itemize}

The benchmark is intentionally transparent, but currently simulated. This boundary is central: the reported values validate the framework's reproducibility and diagnostic behavior, while live models, tools, and organizational workloads are required to establish external validity.

\section{Related Work}

\textbf{Reasoning and tool use.} ReAct interleaves reasoning and environment actions, enabling plans to update from observations~\cite{yao2023react}. Toolformer learns whether, when, and how to call external APIs~\cite{schick2023toolformer}, while ToolLLM scales tool learning to thousands of APIs~\cite{qin2024toolllm}. AnyTool combines hierarchical retrieval, solving, and self-reflection~\cite{du2024anytool}. These systems primarily improve an agent's ability to use tools; our focus is the evaluation of the policy choosing among several orchestration operations.

\textbf{Routing and efficiency.} FrugalGPT and RouteLLM route requests among models to balance quality and cost~\cite{chen2024frugalgpt,ong2025routellm}. LLMCompiler plans parallel function calls and reports improvements in latency, cost, and accuracy~\cite{kim2024llmcompiler}. MetaRoute-Bench extends the routing question across decomposition, tool use, code execution, delegation, verification, and direct response, while retaining raw operational metrics.

\textbf{Agent evaluation.} AgentBench evaluates agents across interactive environments and identifies long-term reasoning and decision-making failures~\cite{liu2024agentbench}. SWE-bench grounds coding-agent evaluation in real repository issues~\cite{jimenez2024swebench}; $\tau$-bench evaluates tool-using conversational agents under domain policies~\cite{yao2024taubench}. These works motivate executable and realistic evaluation. Our current synthetic benchmark is complementary as a controlled diagnostic, but it does not replace live-environment evidence.

\section{Framework}

\subsection{Meta-Decision Process}

Each task $x$ has a workload, difficulty, ambiguity, operation-need annotations, and cost and latency budgets. The present policies map task metadata to a complete route $r$:

\begin{equation}
r = \pi(x), \quad r \subseteq \{D,T,C,G,V,A\},
\end{equation}

where $D$ denotes decomposition, $T$ tool use, $C$ code execution, $G$ delegation, $V$ verification, and $A$ final answering. Every route must end in $A$; operations are unique in the current implementation. A trace records actions, success, expected success, cost, latency, retries, budget compliance, confidence, and failure mode. The framework can be extended to policies over execution history, but the evaluated policy adapts only through one executor-level retry after a failed operation; it does not replan from arbitrary intermediate outputs.

We report success, cost, and latency separately. For ranking in the command-line summary only, we define a secondary utility

\begin{equation}
U = S - 0.025K - 0.0015L,
\end{equation}

where $S$ is success rate, $K$ mean normalized cost, and $L$ mean latency in seconds. Conclusions do not depend solely on this weighting.

\subsection{System Architecture}

The implementation separates five interfaces: (1) a deterministic workload generator; (2) routing policies; (3) a seeded offline executor; (4) typed trace export; and (5) aggregate and paired evaluation. This separation permits replacement of the simulator with live model and tool adapters without changing policy or analysis interfaces.

The executor makes its assumptions explicit. Each operation has a normalized cost, latency, benefit scaled by task need, and---for external tools, code, and delegation---a failure probability. Difficulty and ambiguity reduce base success. Unnecessary operations impose overhead. Adaptive recovery makes one retry after an execution failure. All randomness is keyed by seed, task identifier, and policy name.

\subsection{Policies}

We evaluate four fixed policies (direct, always decompose, always tool, and always code), a random route policy, a workload-specific static rule table, a one-shot router selecting the highest annotated operation need, and the proposed task-aware compositional policy, labeled \emph{adaptive} in the artifacts. It selects up to three operations above a difficulty-dependent threshold, orders decomposition first and verification last, and enables one recovery attempt. Both one-shot and adaptive policies receive the same task-level need annotations; their comparison therefore isolates route composition and recovery rather than feature prediction.

\section{Evaluation}

\subsection{Workloads and Protocol}

The suite contains 60 task profiles for each of data analysis, research, and document processing. Profiles cover four difficulty levels and use seeded variation around workload-specific needs. Data analysis emphasizes code and verification, research emphasizes decomposition and tools, and document processing emphasizes tools and verification. These profiles represent workflow characteristics rather than natural-language task instances.

We execute every policy on every profile for 30 paired seeds. This produces 5,400 traces per policy and 43,200 main-experiment traces. We report mean success and a 95\% confidence interval computed over seed-level success rates, plus mean normalized cost, latency, and cost per success. Paired differences use the same seed-level aggregation. The complete trace file is exported as CSV.

\subsection{Main Results}

\begin{table}[t]
\caption{Main benchmark results. CI is the 95\% interval half-width for success over 30 seeds.}
\label{tab:main}
\small
\begin{tabular}{lrrrr}
\toprule
Policy & Success & CI & Cost & Lat. (s) \\
\midrule
Adaptive & \textbf{.794} & .012 & 2.91 & 20.45 \\
Static workload & .767 & .013 & 2.78 & 19.23 \\
One-shot & .674 & .012 & 1.83 & 12.46 \\
Always tool & .637 & .011 & 1.96 & 13.44 \\
Random & .635 & .012 & 2.54 & 17.86 \\
Always decompose & .610 & .013 & 1.34 & 8.40 \\
Always code & .608 & .011 & 2.24 & 15.68 \\
Direct & .529 & .012 & .73 & 4.48 \\
\bottomrule
\end{tabular}
\end{table}

Table~\ref{tab:main} shows a clear operational tradeoff. Adaptive routing has the highest success rate, but direct answering is least expensive and fastest. Against the strongest baseline, static workload routing, adaptive routing improves success by 2.74 points (paired 95\% CI $\pm$1.96), while increasing cost by 0.13 units and latency by 1.22 seconds. Compared with one-shot routing, adaptive routing improves success by 12.02 points ($\pm$1.87), with 1.09 additional cost units and 8.00 additional seconds.

Workload-level adaptive success is .802 for data analysis, .835 for research, and .746 for document processing. The static policy is particularly competitive for data analysis (.791), but trails more on research (.786 versus .835). This suggests that task-level composition has the most value where decomposition and information access interact.

\subsection{Ablations}

\begin{table}[t]
\caption{Adaptive-policy ablations over 5,400 traces each.}
\label{tab:ablation}
\small
\begin{tabular}{lrrr}
\toprule
Policy & Success & Cost & Lat. (s) \\
\midrule
Adaptive & \textbf{.794} & 2.91 & 20.45 \\
No recovery & .783 & 2.86 & 19.98 \\
No decomposition & .769 & 2.88 & 20.44 \\
No verification & .747 & 2.69 & 18.96 \\
Single operation & .681 & 1.85 & 12.68 \\
\bottomrule
\end{tabular}
\end{table}

Restricting routes to one operation produces the largest reduction: 11.30 points ($\pm$1.85). Removing verification reduces success by 4.70 points ($\pm$1.68), and removing decomposition reduces it by 2.57 points ($\pm$1.68). The 1.13-point recovery difference has an interval of $\pm$1.70 and therefore does not support a confident recovery benefit in this experiment. This is a useful negative finding: recovery is operationally plausible, but the current failure frequency and sample design do not establish its effect.

\section{Operational Lessons}

First, a strong static policy deserves inclusion in agent evaluations. It captures much of the benefit of adaptive routing and is simpler to inspect and operate. Second, success improvements must be reported alongside cost and latency: the adaptive policy is neither cheapest nor fastest. Third, route composition matters more than any single fixed action in this model. Fourth, structured traces make policy behavior auditable; aggregate accuracy alone cannot reveal unnecessary calls, retry behavior, or workload-specific regressions.

For deployment, the framework should sit above existing model and tool adapters. Decisions and observations should be logged without private chain-of-thought, using structured action labels, confidence, tool outcomes, and stop reasons. Production rollout should begin in shadow mode, compare against the incumbent static policy, enforce hard budgets, and support human escalation.

\section{Limitations, Ethics, and Next Steps}

The primary limitation is external validity. Outcomes are sampled from a hand-specified execution model whose operation benefits depend on the same need dimensions supplied to task-aware policies. The adaptive policy should therefore be interpreted as an annotated upper-bound controller, not a learned router. Cost units and latency values are normalized assumptions, not invoices or wall-clock measurements. Task profiles do not contain natural-language inputs, live APIs, concurrent agents, security constraints, or human judgments.

These limitations prevent claims of deployment readiness. The next evaluation must replace task annotations with predictions derived from task text, replay real anonymized traces, and then run controlled live tools under fixed budgets. A learned router, model-routing baselines, calibration analysis, and human escalation should also be added. We publish all simulator constants and raw traces so these assumptions can be challenged rather than hidden.

Agent routing can create privacy, security, and accountability risks when tasks are delegated to external services or code is executed. Practical deployments need data-minimization rules, sandboxed execution, allowlisted tools, access control, audit logs, and human review for high-impact actions. The current benchmark performs no real external action and contains no personal or proprietary data.

\section{Conclusion}

MetaRoute-Bench isolates the orchestration policy governing how an agentic system decomposes, routes, executes, verifies, and recovers. In a reproducible offline study, task-aware adaptive composition improves success over fixed and one-shot policies while incurring measurable cost and latency. The framework and traces provide a foundation for the more consequential next step: validation on live, production-representative workflows.

\bibliographystyle{ACM-Reference-Format}
\bibliography{references}

\end{document}
